% =============================================================
% BILAGA: Fördefinierade träningsmallar
% Översikt av de mallar som används vid regelbaserad fallback
% och som komplement till den AI-baserade träningsplaneringen.
% =============================================================

\section{Fördefinierade träningsmallar}
\label{app:mallar}

Denna bilaga ger en överblick av de fördefinierade träningsmallar som används
i prototypen. Mallarna utgör grunden för den regelbaserade fallbacken som
aktiveras när RAG-lösningen inte kan användas eller när den AI-genererade
planen inte kan repareras till regelkompatibilitet. Den fullständiga
implementationen finns i projektets kodrepository~\cite{Boissier_AI-baserat_traningsstod_for}.

Mallstrukturen bygger på tre nivåer: en \emph{veckomall} som anger
svårighetsgrad per dag, ett \emph{passbibliotek} med konkreta passförslag
indexerade på disciplin, träningsfas och svårighet, samt en \emph{disciplinmappning}
som styr vilken kategori av mallar som används för olika grenar.

% ----------------------------------------------------------------
\subsection{Veckomallar}
\label{app:mallar:vecka}

Beroende på hur många träningsdagar idrottaren har per vecka väljs en
fördefinierad veckostruktur. Varje dag tilldelas en svårighetsgrad:
\textit{hard} (kvalitetspass), \textit{medium} (teknik eller medelintensiv
löpning), \textit{easy} (lugn distans eller återhämtning) eller \textit{rest}
(vila). Tabell~\ref{tab:veckomallar} visar de fyra veckomallarna.

\begin{table}[h]
\centering
\caption{Veckomallar per antal träningsdagar.}
\label{tab:veckomallar}
\small
\begin{tabular}{l c c c c c c c}
\hline
\textbf{Dagar/v} & \textbf{Mån} & \textbf{Tis} & \textbf{Ons} & \textbf{Tor} & \textbf{Fre} & \textbf{Lör} & \textbf{Sön} \\
\hline
3 & hard   & rest   & medium & rest & hard & rest   & rest \\
4 & hard   & easy   & medium & rest & hard & easy   & rest \\
5 & hard   & easy   & medium & easy & hard & medium & rest \\
6 & hard   & easy   & medium & easy & hard & medium & rest \\
\hline
\end{tabular}
\end{table}

Strukturen följer principen om hård-lätt-växling, där högintensiva pass
varvas med återhämtning eller lugnare pass.

% ----------------------------------------------------------------
\subsection{Passbibliotek}
\label{app:mallar:pass}

Varje passmall innehåller fem fält: \textit{namn}, \textit{typ} (snabbhet,
uthållighet, teknik eller återhämtning), \textit{längd} i minuter,
\textit{intensitet} (hög, medel eller låg) samt en detaljerad
\textit{beskrivning} med uppvärmning, huvudpass och nedjogg. Mallarna är
organiserade i en tredimensionell struktur: disciplin $\times$ träningsfas
$\times$ svårighet.

\begin{itemize}
    \item \textbf{Discipliner:} sprint (100--200~m), medel (400--1500~m) och
    distans (3000~m och längre).
    \item \textbf{Träningsfaser:} grundträning, uppbyggnad, tävling och
    återhämtning.
    \item \textbf{Svårigheter:} hard, medium och easy.
\end{itemize}

Totalt ger detta $3 \times 4 \times 3 = 36$ kategorier av mallar, där varje
kategori innehåller mellan ett och två passförslag.
Tabell~\ref{tab:passbibliotek_sprint}, \ref{tab:passbibliotek_medel} och
\ref{tab:passbibliotek_distans} visar exempel på pass per disciplin och fas
(svårighet \textit{hard} och \textit{easy}).

\begin{table}[h]
\centering
\caption{Exempel på passmallar för sprint (100--200~m).}
\label{tab:passbibliotek_sprint}
\small
\begin{tabular}{l l p{0.45\textwidth}}
\hline
\textbf{Fas} & \textbf{Svårighet} & \textbf{Pass} \\
\hline
Grundträning & hard & Accelerationsträning (60 min), Fartträning 150~m \\
Grundträning & easy & Lugn löpning (35 min), Återhämtning \& rörlighet \\
Uppbyggnad   & hard & Maxfart 60~m, Fartuthållighet sprint \\
Uppbyggnad   & easy & Lätt jogg (30 min), Aktiv vila \\
Tävling      & hard & Tävlingsförberedelse (3$\times$30~m + 2$\times$60~m) \\
Tävling      & easy & Lätt löpning (25 min), Vila + stretch \\
Återhämtning & easy & Lätt jogg (25 min), Stretch \& rörlighet \\
\hline
\end{tabular}
\end{table}

\begin{table}[h]
\centering
\caption{Exempel på passmallar för medeldistans (400--1500~m).}
\label{tab:passbibliotek_medel}
\small
\begin{tabular}{l l p{0.45\textwidth}}
\hline
\textbf{Fas} & \textbf{Svårighet} & \textbf{Pass} \\
\hline
Grundträning & hard & Tröskelpass (60 min), Intervaller 8$\times$400~m \\
Grundträning & easy & Lugn distanslöpning (45 min), Återhämtningsjogg \\
Uppbyggnad   & hard & VO2max 5$\times$1000~m, Tempointervaller 6$\times$600~m \\
Uppbyggnad   & easy & Lugn löpning (40 min), Aktiv vila \\
Tävling      & hard & Tävlingsspecifikt (3$\times$300~m + 2$\times$200~m) \\
Tävling      & easy & Lätt löpning (30 min), Vila + förberedelse \\
Återhämtning & easy & Lugn jogg (30 min), Stretch \& rörlighet \\
\hline
\end{tabular}
\end{table}

\begin{table}[h]
\centering
\caption{Exempel på passmallar för distans (3000~m och längre).}
\label{tab:passbibliotek_distans}
\small
\begin{tabular}{l l p{0.45\textwidth}}
\hline
\textbf{Fas} & \textbf{Svårighet} & \textbf{Pass} \\
\hline
Grundträning & hard & Tröskelpass långt (2$\times$15 min), Intervaller 6$\times$1000~m \\
Grundträning & easy & Lång löpning (60 min), Återhämtningsjogg \\
Uppbyggnad   & hard & VO2max 5$\times$1200~m, Tempolöpning (25 min tröskel) \\
Uppbyggnad   & easy & Lugn distanslöpning (50 min), Aktiv vila \\
Tävling      & hard & Tävlingsspecifikt 3$\times$1000~m \\
Tävling      & easy & Lätt löpning (30 min), Vila + stretch \\
Återhämtning & easy & Mycket lugn jogg (30 min), Stretch \& rörlighet \\
\hline
\end{tabular}
\end{table}

Vid val av pass kombineras dagens svårighet i veckomallen med idrottarens
disciplin och aktuella träningsfas, varefter ett pass slumpas ur den
matchande listan.

% ----------------------------------------------------------------
\subsection{Disciplinmappning}
\label{app:mallar:mappning}

Eftersom mallbiblioteket är uppbyggt kring de tre löpkategorierna sprint,
medel och distans används en mappning för grenar utanför dessa, vilket
visas i tabell~\ref{tab:disciplinmappning}.

\begin{table}[h]
\centering
\caption{Mappning från friidrottsgren till löpkategori i mallbiblioteket.}
\label{tab:disciplinmappning}
\small
\begin{tabular}{l l}
\hline
\textbf{Gren} & \textbf{Använder mallar för} \\
\hline
Sprint    & sprint \\
Medel     & medel \\
Distans   & distans \\
Hopp      & sprint \\
Kast      & sprint \\
Mångkamp  & medel \\
\hline
\end{tabular}
\end{table}

Mappningen utgår från att hoppare och kastare ofta tränar med
sprintliknande inslag, medan mångkampare har en blandad träningsprofil
som närmast motsvarar medeldistans.

% ----------------------------------------------------------------
\subsection{Exempel på passmall}
\label{app:mallar:exempel}

Som illustration visas innehållet i en enskild passmall, här
\textit{Tröskelpass} för medeldistans i grundträningsfas:

\begin{quote}
\textbf{Namn:} Tröskelpass \\
\textbf{Typ:} Uthållighet \\
\textbf{Längd:} 60 min \\
\textbf{Intensitet:} Hög \\
\textbf{Beskrivning:} \\
Uppvärmning: 15 min jogg + dynamisk stretch \\
Huvudpass: 20 min löpning i tröskeltempo (puls $\sim$85\,\% av max) \\
Nedjogg: 15 min + stretch
\end{quote}

Övriga mallar följer samma struktur och kan inspekteras i sin helhet
i projektets kodrepository~\cite{Boissier_AI-baserat_traningsstod_for}.
